%===================================================================
% Chapter 6: Related Work
% Consolidated thematic survey across the three systems.
%===================================================================
\chapter{Related Work}
\label{chap:related}

The three systems of this dissertation draw on, and depart from,
several bodies of prior work. Rather than surveying them
system-by-system, this chapter organizes the discussion by theme:
rule-based query rewriting, the departure point for the entire
dissertation (\cref{sec:related:rules}); the small body of work that,
like Chapters 3 and 4, rewrites queries without being confined to
rules (\cref{sec:related:beyondrules}); the techniques for
establishing query equivalence that every rewriter in this
dissertation depends on (\cref{sec:related:equiv}); the broader uses
of program synthesis in data systems, and the provenance literature
adjacent to Chapter 3's dataflow abstraction
(\cref{sec:related:synthesis}); the rapidly growing application of
large language models across the database stack
(\cref{sec:related:llms}); and the techniques that, like Chapter 5,
look beyond a single query (\cref{sec:related:pipelines}).

\section{Rule-Based Query Rewriting}
\label{sec:related:rules}
\label{slab:sec:related}

The dominant tradition in query rewriting encodes optimization
knowledge as rewrite rules, and its history is one of steadily
widening sources for those rules. The classical source is human
expertise: rule sets crafted and refined by database experts over
decades~\cite{dintyala2020sqlcheck,lohman1988grammar,finance1991rule,pirahesh1992extensible,chen2016memsql,begoli2018apache,graefe1987rule,ahmed2006cost,levy1994query,muralikrishna1992improved,seshadri1996cost}.
Such collections grow slowly, and each rule must anticipate a pattern
and encode its correctness conditions by hand. Automatic rule
discovery lightens the authoring burden: \wetune~\cite{wetune}
enumerates and verifies new rules mechanically, though over a
restricted space of simple query shapes, and subsequent systems
propose rules tailored to analytical
workloads~\cite{sarthi2020generalized,blitz2,modi2021new,bruno2022computation}.
Orthogonally, learned techniques improve how a fixed rule set is
applied: \learnedrewrite~\cite{LR} uses Monte Carlo tree search with
learned cost models to choose the order in which Calcite's rules fire.
A separate line exploits integrity constraints through rules, using
key and \texttt{NOT NULL} constraints to eliminate
joins~\cite{join-elimination}, optimize queries with
disjunction~\cite{blog-disjunct}, or remove unnecessary
\sqldistinctcolor~\cite{habimana2015query}, among
others~\cite{presto-qo,db2-qo}; these constraint-driven rewrites
remain pattern-matched and limited to relatively simple constraint
forms, where the synthesis of \cref{slab:chap} can exploit richer
integrity constraints freely during search. Most recently, large
language models have been enlisted in service of rules:
\rbot~\cite{sun2024r} combines retrieval and self-reflection to apply
rules well, and \textsc{LLM-R\textsuperscript{2}}~\cite{li2024llm}
uses contrastive demonstrations to select effective rules from a given
set.

However their rules are obtained, ordered, or selected, all of these
systems share the boundary this dissertation set out from: they cannot
optimize a query whose opportunity matches none of their rules.
\Cref{slab:chap} demonstrated this boundary empirically, on queries no
state-of-the-art rule-based system could improve.

\section{Query Rewriting Beyond Rules}
\label{sec:related:beyondrules}
\label{gen:sec:related}

Before this dissertation, two systems had applied program synthesis to
source-to-source query rewriting, and each kept rules or locality in
the loop. FGH-rule~\cite{wang2022optimizing} operates on Datalog: a
rule matches a recursion pattern in the input query and produces an
output template with an unknown expression, which synthesis then fills
so that the completed query is equivalent to the input. The rule
delimits the synthesis search space, so the approach's reach is the
rule's reach, and its formal guarantees have been demonstrated at
small scale (seven queries). \sia~\cite{zhou2021sia} uses synthesis to
add predicates to a query, but only within the \sqlwherecolor clause;
it cannot restructure the rest of the query. \slabcity
(\cref{slab:chap}) is, to the best of our knowledge, the first
synthesis-based rewriter that is capable of whole-query optimization
and needs no rewrite rules at all.

The same ambition, rewriting beyond what rules capture, has since been
pursued with language models used directly rather than in service of a
rule set. LITHE~\cite{llm-rw2} prompts the model with one of six
handcrafted natural-language rules, each accompanied by an example;
the rule set is fixed, so the system's knowledge cannot grow with
database feedback. QUITE~\cite{song2025quite} casts rewriting as a
Markov decision process whose refinement actions are drawn from a
curated knowledge base of documented heuristics, which confines the
search to known actions. \genrewrite (\cref{gen:chap}) differs from
both in kind: its Natural Language Rewrite Rules are not fixed or
curated in advance but distilled by the model from its own successful
rewrites, so the repository grows over the life of a workload and can
contain strategies present in no existing literature. These systems,
along with \rbot and \textsc{LLM-R\textsuperscript{2}} above, were
published after \genrewrite's initial arXiv
release~\cite{genrewrite-report}.

\section{Checking Query Equivalence}
\label{sec:related:equiv}

Every rewriter in this dissertation admits a candidate rewrite only
after evidence of equivalence, and the available evidence spans the
spectrum Chapter~2 laid out. On the verification end,
\cosette~\cite{chu2017cosette} constructs mechanized proofs of
equivalence with a proof assistant, or counterexamples for
inequivalent pairs; EQUITAS and
\spes~\cite{zhou2019automated,zhou2022spes} execute queries
symbolically and discharge the resulting formulas with an SMT solver,
covering select-project-join queries and some outer joins and
aggregates. \Cref{slab:chap} consumes these verifiers as its final
checking layers, and their partial language coverage is exactly why
its checker is layered rather than monolithic.

On the testing end, mutation-based techniques from the software
engineering literature generate databases to kill query
mutants~\cite{chandra2015data,shah2011generating}, driven by
hard-coded constraint rules or specifications, over query languages
too small for the workloads in this dissertation.
EvoSQL~\cite{castelein2018search} searches for test databases offline,
guided by a predicate coverage metric~\cite{tuya2010full}, at a cost
of minutes per query that no synthesis loop can afford.
RATest~\cite{miao2019explaining} finds minimal distinguishing
databases via a provenance-based algorithm, but is expensive and
supports a limited language. The tester of \cref{slab:chap} differs
from all of these in being domain-specialized and online: it mines the
candidate and input queries for the specific patterns that commonly
make rewrites wrong, and turns those patterns into distinguishing
databases fast enough to sit inside a CEGIS loop. Most recently,
LLM-SQL-Solver~\cite{zhao2023llm} asks a language model to judge
equivalence directly, under a pragmatically relaxed notion of
equivalence; such judgments are useful signals but carry no guarantee,
which is why \cref{gen:chap} treats the model's own equivalence
opinion as a starting point to be checked, not a verdict.

\section{Program Synthesis and Provenance in Data Systems}
\label{sec:related:synthesis}

Beyond rewriting, program synthesis has served data systems in several
roles. Blitz~\cite{synthesis-spark,blitz2} synthesizes user-defined
operators from Spark programs to improve parallel execution;
Chestnut~\cite{chestnut} synthesizes a data layout together with a
query plan over that layout; both target non-relational settings
rather than optimizing core SQL queries.
\textsc{Sickle}~\cite{zhou2022synthesizing} and
Scythe~\cite{wang2017synthesizing} synthesize analytical queries from
input-output examples, easing the user's burden rather than the
query's cost. Synthesis has also been used to generate diverse test
databases for database
applications~\cite{qex,tesma,syndb,xie:symbolic2}, where the goal is
application code coverage.

A reader may also wonder how the dataflow abstraction of
\cref{slab:chap} relates to data provenance, a concept it
superficially resembles. Provenance is studied across probabilistic
databases~\cite{green2006models,suciu2011probabilistic}, view
maintenance~\cite{buneman2002propagation}, and the explanation of
query results~\cite{green2007provenance,chapman2009not}, and has more
recently guided synthesis from
examples~\cite{raghothaman2020provenance,zhou2022synthesizing}. The
difference is granularity and purpose: provenance accounts for where
every individual piece of data originates and how it is computed,
while a dataflow records only the column-level movement of information
through a query's operations. That coarseness is the point: dataflows
are cheap enough to extract and compare inside a search loop, yet
informative enough to steer it.

\section{Large Language Models in Databases}
\label{sec:related:llms}

The systems of Chapters 4 and 5 sit within a rapid broadening of LLM
applications across the database stack. Language models now drive
NL2SQL with strong benchmark results~\cite{li2023can}; infer column
correspondences for schema matching~\cite{parciak2407schema};
synthesize data for augmentation~\cite{tian2025text}; tune and
diagnose systems, as in DB-BERT and
D-Bot~\cite{trummer2022db,zhou2023d}; and integrate more broadly with
DBMS internals for optimization and indexing in
DB-GPT~\cite{zhou2023dbgpt}. Collectively, these works position LLMs
as versatile assistants across the stack. What distinguishes this
dissertation's use of them is the discipline built around the model:
\cref{gen:chap} surrounds it with accumulated rewrite knowledge,
correction loops, and verification, and \cref{dag:chap} with
structural analysis, separated criticism, and data-backed checking, so
that in both cases the model proposes and the system disposes.

\section{Optimizing Beyond the Single Query}
\label{sec:related:pipelines}
\label{dag:sec:related}

\Cref{dag:chap} is, to the best of our knowledge, the first system to
treat an explicit pipeline DAG as the unit of dependency-aware
rewriting, restructuring node boundaries and deciding refactoring,
materialization, and refresh frequency jointly. Four established lines
each cover part of that territory. Single-query rewriting, in all
three flavors this dissertation has discussed, and including
Chapters 3 and 4 themselves, transforms one statement and leaves the
surrounding dependency graph untouched. Multi-query
optimization~\cite{sellis1988multiple,roy2000efficient} and
common-subexpression sharing~\cite{scopeCSE,similarSubexpr} reuse
intermediate results across a batch of statements, and
materialized-view
selection~\cite{agrawal2000automated,harinarayan1996implementing}
precomputes chosen results and redirects consumers to them; these
techniques match on syntactic structure and choose persistence among
results that already exist, without recognizing the same computation
written in dissimilar SQL or trimming results no downstream output
needs. Incremental view maintenance~\cite{ivm} and streaming
materialized-view systems~\cite{materialize,risingwave} cheapen each
refresh by processing only the rows that change; \cref{dag:chap}
approaches freshness from the opposite direction, aligning how often
each node runs with how often its inputs change and its outputs are
consumed, and the two directions are complementary and could compose.
Finally, workflow orchestrators, including dbt~\cite{dbtWhatExactly},
SQLMesh~\cite{sqlmeshOverview}, Airflow~\cite{airflowDags},
Dagster~\cite{dagsterAssets}, Prefect~\cite{prefectTasks},
Argo~\cite{argoDag}, Luigi~\cite{luigiDocs}, and Databricks
Workflows~\cite{databricksJobs}, express exactly the DAGs
\cref{dag:chap} targets, but their schedulers optimize execution
order, invalidation, retries, and resource allocation while leaving
each node's SQL semantics untouched, which is where \cref{dag:chap}'s
reductions originate.
